
\chapter{עבודות קשורות: תיאום ומחקר סוכנים}

\section{סוכנים עם כלים: הדור הראשון}

\textenglish{MRKL} (\textenglish{Modular Reasoning, Knowledge and Language}) (2022) היה ממבשרי הגישה: מנוע ניתוב מבוסס כללים הכיל \textenglish{LLM} ל\textenglish{routing} בין מודולים מיוחדים (חישוב, חיפוש, מסד נתונים). \textenglish{ReAct} \cite{yao2022react} פישט זאת לפרדיגמה אחת גמישה.

\section{ממשק כלים}

\textenglish{Toolformer} \cite{schick2023toolformer} הראה שמודל \textenglish{GPT-J} בגודל \textenglish{6B} יכול ללמוד שימוש ב-5 כלים (\textenglish{Wikipedia, calculator, calendar, translator, QA}) ולהשתוות ל-\textenglish{GPT-3 175B} על מטלות ספציפיות --- רק עם שימוש חכם בכלים. \textenglish{ToolLLM} \cite{qin2023toolllm} הרחיב זאת לאלפי כלים, עם \textenglish{fine-tuning} על \textenglish{ToolBench}.

\section{מסגרות רב-סוכן}

\textenglish{AutoGen} \cite{wu2023autogen} ו-\textenglish{MetaGPT} \cite{hong2023metagpt} הדגימו שחלוקת תפקידים בין סוכנים מפחיתה שגיאות ומשפרת מורכבות עבור פרויקטים גדולים. \textenglish{MetaGPT} הציג ביצועים מובילים על \textenglish{HumanEval} ו-\textenglish{MBPP} לכתיבת קוד, בעיקר בזכות שלב ה-\textenglish{design document} שמבנה את הפעולה.

\section{מחקרי עלות ויעילות}

\cite{segal2024orchestration} הגדיר את נוסחת \textenglish{WC} לעלות אסימון ויישם ניתוח על \textenglish{CrewAI}. תוצאותיהם מראות כי הפחתת \textenglish{context window} ב-\textenglish{30\%} (על ידי ניהול חכם של היסטוריה) מפחיתה עלות ב-\textenglish{23\%} ללא פגיעה בביצועים. גישת ה-\textenglish{Router-Skill} שלנו משלימה ממצא זה על ידי הפחתת אסימוני טעינת כישורים.

\section{כיוון משולב}

\cite{nakash2024skillsieve} מציג ראיות ש\textenglish{SkillSieve} ניתן לשילוב עם כל אחת מהמסגרות הנ"ל (ReAct, AutoGen, CrewAI) ללא שינוי בממשק הסוכן, כ-\textenglish{middleware} שקוף.
